\documentclass[11pt,a4paper]{article}

\usepackage[utf8]{inputenc}
\usepackage[T1]{fontenc}
\usepackage[margin=2.4cm]{geometry}
\usepackage{amsmath,amssymb}
\usepackage{booktabs}
\usepackage{graphicx}
\usepackage{xcolor}
\usepackage{hyperref}
\hypersetup{colorlinks=true,linkcolor=blue!55!black,urlcolor=blue!55!black,
            citecolor=blue!55!black}
\usepackage{caption}
\captionsetup{font=small,labelfont=bf}

\graphicspath{{figures/}}

\newcommand{\Prob}{\mathbb{P}}
\newcommand{\Ex}{\mathbb{E}}
\newcommand{\KL}[2]{D\!\left(#1 \,\|\, #2\right)}
\newcommand{\A}[1]{\texttt{ALLOCATION\_#1}}

\title{\textbf{Pre-analysis of \texttt{task\_003}: the task, the controller
defect, and a coarse-graining that makes information quantities estimable}}
\author{MA-CC blackboard-control study}
\date{2 October 2026}

\begin{document}
\maketitle

\begin{abstract}
\noindent
This report consolidates an offline re-analysis of the first batch of
blackboard-control simulations (study \texttt{21-09-2026-full-vs-report-v1}, 60
initializations of \texttt{task\_003} with \texttt{gpt-oss-120b}). It reaches
three conclusions. First, \texttt{task\_003} is fully reconstructible from its
nine hidden latent values, and the population of agents collectively holds the
only short proof of the correct answer. Second, the controller used in that
batch is \emph{not} what its configuration claims: fact admission was evaluated
once for the decoy target and reused for both steering directions, so the
truth-targeting arm recommends \A{0} while quoting evidence selected to
undermine it. Measured target information is $0.022$ bits --- essentially zero.
Third, the raw population state (a 325-valued vote vector) defeats every
estimator we tried; binning the mean agent posterior into four states makes the
path divergence \emph{exactly} computable, but four independent tests agree the
resulting chain is Markovian only in the silent arm at persistence $\rho = 1$.
No simulations were run and no provider requests made.
\end{abstract}

\tableofcontents

\section{Scope and vocabulary}

Three terms recur and are defined once here.

\paragraph{Blackboard.} A shared message board. Each round, every agent may post
a statement; all posts are visible to all agents in that round, and the board is
cleared before the next round. Nothing persists on the board itself.

\paragraph{Epistemic persistence $\rho$.} Each agent remembers each fact it holds
with probability $\rho$, independently per agent, per fact, per round. At
$\rho = 1$ memory is perfect; at $\rho = 0.75$ a quarter of each agent's
knowledge decays every round and must be re-acquired from the board to survive.

\paragraph{Controller.} An extra participant that may post selected true facts to
steer the population toward a named allocation (its \emph{target}). It is
\emph{silent} in the baseline arm. In the intervention arms its target is either
\A{0} (the correct answer, called the \emph{truth} direction) or \A{2} (an
incorrect answer, the \emph{decoy} direction).

\section{What \texttt{task\_003} is}

\subsection{The underlying problem}

The task is a Team Allocation problem in the style of MuSR. Three people must be
split between two jobs: one person works alone on the first job, the remaining
two work jointly on the second. There are therefore exactly three candidate
allocations, one per choice of the lone worker.

Nine hidden integers in $\{1,2,3\}$ describe the world: for each of the three
people, a skill level at each of the two jobs (six numbers), plus a cooperation
level for each of the three possible pairs (three numbers). An allocation's score
is
\begin{equation}
  \text{score} = \underbrace{s_{\text{singleton},1}}_{\text{lone worker's skill at job 1}}
  \;+\; \underbrace{\sum_{p \,\in\, \text{pair}} s_{p,2}}_{\text{pair's skills at job 2}}
  \;+\; \underbrace{c_{\text{pair}}}_{\text{how well the pair cooperates}} .
\end{equation}
The highest-scoring allocation is the correct answer.

\subsection{The specific world}

\texttt{task\_003} is the latent vector
\begin{equation}
  V = (3,1,1,2,2,1,\;1,1,2),
\end{equation}
reading the six skills first (person by person, job by job) then the three
cooperation values. Scoring gives \A{0} as the unique winner. We rebuilt the task
locally from $V$ alone and verified the result is bit-identical to the archived
frozen task: the cluster copies are not needed to reproduce anything in this
report.

\subsection{Facts}

A \emph{fact} is a true natural-language proposition about the world, drawn from
a fixed catalog of 99 propositions: 45 \emph{threshold} statements (``person $p$
has at least moderate skill at job $j$'') and 54 \emph{comparison} statements
(``$p$ is better than $p'$ at job $j$''). Of these, \textbf{49 are true in
\texttt{task\_003}} and only those can ever be posted. Nobody lies in these
simulations; steering is done purely by \emph{which} true facts get said.

\subsection{Exact posterior}

Because the world is only nine ternary numbers, we enumerate all
$3^9 = 19{,}683$ candidate worlds, keep the $14{,}388$ with a unique winner, and
for any set $S$ of facts compute the posterior exactly by filtering:
\begin{equation}
  \Prob(\A{k} \mid S) = \frac{\#\{w : w \models S,\; \text{winner}(w) = k\}}
                            {\#\{w : w \models S\}} .
  \label{eq:posterior}
\end{equation}
This is a count, not an estimate. It is the ground truth against which every
agent's belief is scored. Implementation: \texttt{10\_task\_and\_facts/engine.py}.

\subsection{Who can prove the answer}

A set $S$ \emph{solves} the task if every surviving world has \A{0} as winner,
i.e.\ $\Prob(\A{0} \mid S) = 1$. Enumerating minimal solving sets:

\begin{center}
\begin{tabular}{lr}
\toprule
quantity & value \\
\midrule
smallest solving set & \textbf{4 facts} \\
number of 4-fact proofs & 24 \\
minimal proofs up to size 6 & 7{,}694 \\
minimal proofs the agents can jointly assemble & \textbf{4} \\
minimal proofs the controller's pool can assemble & \textbf{0} \\
facts held by nobody & 12 \\
\bottomrule
\end{tabular}
\end{center}

\noindent
The fact \texttt{cf\_x00\_eq\_3} appears in \emph{every} 4-fact proof and is in
nobody's starting hand. The population therefore holds the only short route to
certainty, and barely: four proofs out of 7{,}694.

\paragraph{Private-eligibility.} A fact is eligible to be handed out privately
when it is individually weak: $\max_k \Prob(\A{k} \mid \{f\}) \le 0.45$
\emph{and} its normalized posterior entropy is $\ge 0.90$. Both conditions must
hold. 41 of the 49 facts qualify; the remaining 8 are too revealing to
distribute and are given to no one.

\paragraph{The decisive set.} A greedy search over the 41 private-eligible facts
produces a 6-fact set reaching $\Prob(\A{0}) = 1$. These six are the
\emph{decisive} facts. They are a \emph{subset} of the private-eligible ones ---
the two categories are nested, not disjoint, which is easy to double-count.

\section{The controller defect}

\subsection{What went wrong}

Admission of a fact into the controller's pool is decided by whether it moves the
posterior toward the target, $\Delta\Prob(\text{target} \mid f) > 0$, together with
a non-decisiveness condition. In the archived batch this test was evaluated
\textbf{once, for the decoy target \A{2}}, and the resulting 24-fact pool was
reused for both configurations. Excluding decisive facts removes only 3 of them,
all \A{2}-favouring.

The consequence is stark: in the truth arm the controller \emph{recommends} \A{0}
while \emph{citing} facts chosen to push the population toward \A{2}. The two
arms differ mainly in how often the controller fires at all --- 32\% of rounds in
the truth arm against 54\% in the decoy arm --- not in what they say.

\begin{quote}
\textbf{Standing caveat.} The \A{0} arm is confounded. Its numbers must not be
read as evidence about truth-directed steering. Its recorded \emph{dynamics} are
sound and are used wherever only dynamics matter.
\end{quote}

\subsection{The replacement}

A corrected design uses two \emph{matched symmetric pools}: 12 facts each,
disjoint, built by running the admission test separately per target, and
verified so that neither pool can prove its own target (otherwise the controller
would simply hand over the answer). These are stored in
\texttt{10\_task\_and\_facts/symmetric\_pools.json}.

\section{Coarse-graining the population state}

\subsection{Why the raw state fails}

The natural state --- the vector of which allocation each agent currently
favours --- takes 325 values in these runs. With 1{,}800 to 5{,}000 transitions
per cell, most of those states are visited once or never. Critic-based
(variational) estimators of the path divergence on this space failed every
reliability check we applied: the estimates did not stabilize with more data and
their confidence intervals did not cover their own point estimates.

\subsection{The epistemic coordinate}

We replace the vote vector with a single number: the average correctness of
belief across the population,
\begin{equation}
  e_t \;=\; \frac{1}{N}\sum_{j=1}^{N} \Prob\!\left(\A{0} \mid K_j(t)\right),
  \label{eq:ecoord}
\end{equation}
where $K_j(t)$ is the set of facts agent $j$ currently remembers and the
posterior is the exact count of Eq.~\eqref{eq:posterior}. So $e_t$ is not a
vote share and not an opinion average; it is ``how close to the truth the
population's evidence actually puts it,'' measured on the exact posterior.

We then bin $e_t$ into \textbf{four states} with cut points at
$0.33$, $0.50$, $0.75$. The cuts are semantic rather than data-driven: $1/3$ is
the prior (knowing nothing), $0.5$ is a bare majority of surviving worlds, and
$0.75$ is near-certainty. A sweep over 2--6 bins and over quantile versus
semantic cuts (Fig.~\ref{fig:bins}) shows 3, 4 and 5 bins are statistically
indistinguishable once resampled across 12 seeds; 4 was chosen for data budget.
An 8-state version that also tracks knowledge \emph{spread} repairs the
$\rho = 1$ cells but leaves three rows of the transition matrix estimated from
too few transitions, so 4 states are used everywhere.

\begin{quote}
\textbf{Never pool persistence levels.} The binning and the fitted kernel are
estimated separately for each $\rho$ and each arm.
\end{quote}

\subsection{Why this buys exactness}

On a finite state space the divergence between two \emph{trajectory laws} has a
closed form. If $P$ has initial distribution $p_0$ and kernel $T^P$, and $Q$ has
$q_0, T^Q$, then over $h$ steps
\begin{equation}
  \KL{P}{Q} \;=\; \KL{p_0}{q_0}
  \;+\; \sum_{t=0}^{h-1} \sum_{i} \mu_t(i) \, \KL{T^P(\cdot \mid i)}{T^Q(\cdot \mid i)},
  \label{eq:pathkl}
\end{equation}
with $\mu_t$ the occupancy of $P$ at time $t$. Every term is a finite sum over
four states. \textbf{There is no critic and no variational bound} --- that is the
entire point of the coarse-graining.

\section{Is the coarse-grained process Markov?}

Equation~\eqref{eq:pathkl} is exact \emph{for a Markov chain}. Whether the binned
process is one is an empirical question, and we test it four independent ways.

\begin{enumerate}
  \item \textbf{Order comparison.} Fit a first-order and a second-order kernel;
        report the held-out log-likelihood gain per step. A Markov process gains
        nothing.
  \item \textbf{Chapman--Kolmogorov.} Compare the empirical $n$-step transition
        matrix with the $n$-th matrix power of the one-step kernel, in total
        variation. A Markov process agrees at all $n$.
  \item \textbf{Conditional independence.} Estimate
        $I(s_{t+1}; s_{t-1} \mid s_t)$ and compare it with a permutation null.
        A Markov process has zero.
  \item \textbf{Lumpability.} Check whether states that were merged behave
        identically, by splitting each bin on an extra coordinate and measuring
        how far the two halves' outgoing rows differ.
\end{enumerate}

\begin{figure}[t]
  \centering
  \includegraphics[width=\textwidth]{fig_tests}
  \caption{The first three tests, per persistence level and arm. Green marks the
  one cell that passes: silent at $\rho = 1$. Test 1 is negative there (a
  second-order model fits \emph{worse} out of sample), the
  Chapman--Kolmogorov gap stays near $0.03$ out to five steps, and the
  conditional mutual information is within noise of its permutation null
  ($p = 0.15$; every other cell gives $p \approx 0.005$).}
  \label{fig:tests}
\end{figure}

\begin{table}[t]
  \centering
  \small
  \begin{tabular}{llrrrrr}
    \toprule
    $\rho$ & arm & $n$ & order-2 gain & CK gap ($n{=}5$) & CMI excess (bits) & $p$ \\
    \midrule
    0.75 & silent & 1800 & $0.0161$ & $0.106$ & $0.0284$ & 0.005 \\
    0.75 & A2     & 5055 & $0.0146$ & $0.109$ & $0.0261$ & 0.005 \\
    0.75 & A0     & 5205 & $0.0139$ & $0.110$ & $0.0225$ & 0.005 \\
    1.00 & silent & 1800 & $\mathbf{-0.0050}$ & $\mathbf{0.033}$ & $\mathbf{0.0017}$ & \textbf{0.154} \\
    1.00 & A2     & 2520 & $0.0060$ & $0.078$ & $0.0130$ & 0.005 \\
    1.00 & A0     & 4350 & $0.0025$ & $0.068$ & $0.0074$ & 0.005 \\
    \bottomrule
  \end{tabular}
  \caption{All four tests agree. Only $\rho = 1$ / silent is Markovian in the
  four-state coordinate. At $\rho = 0.75$ roughly $0.015$ nats per step of
  memory is unexplained.}
  \label{tab:tests}
\end{table}

\subsection{Why lower persistence makes this harder, not easier}

The intuition that forgetting destroys memory, and so should \emph{help}
Markovianity, is the wrong way round. Forgetting destroys memory in the agents;
it does not destroy memory in the \emph{coarse variable}. Two populations with
the same average correctness $e_t$ can differ in how that correctness is
distributed --- one where every agent knows a little, one where a few agents know
a lot. Under perfect memory that difference is frozen and does not affect the
next step. Under decay the concentrated population loses more, because its
information sits in fewer places, so the two populations move differently from
the same value of $e_t$. The hidden coordinate is knowledge \emph{spread}, and
decay is exactly what makes it matter. The lumpability test confirms it:
splitting on spread gives the largest row discrepancy in the silent arm
($0.130$ at $\rho = 0.75$), while under control the worst offender is instead
the survival of decisive facts.

\section{Information estimates}

\subsection{Quantities}

\paragraph{Path divergence.} $K = \KL{P^{\text{ctl}}}{P^{\text{silent}}}$ from
Eq.~\eqref{eq:pathkl}: how far intervention displaces the trajectory law from
silence. Units: nats.

\paragraph{Intervention information.} $I(\text{intervention}; s_{0:h})$, at most
1 bit: can you tell from the trajectory alone whether a controller was running?

\paragraph{Target information.} $I(\text{target}; s_{0:h})$, also at most 1 bit
given two targets: can you tell \emph{which way} it was steering? This is the
quantity a control-efficiency claim rests on.

\paragraph{Efficiency.} $\eta = I(\text{target};\cdot) / C_\pi$, where $C_\pi$ is
the control cost measured as the accumulated divergence. Bits of achieved
steering per nat of intervention.

\subsection{Results at horizon $h = 14$}

\begin{table}[h]
  \centering
  \small
  \begin{tabular}{lrrrrr}
    \toprule
    $\rho$ & $K(\text{A2} \| \text{silent})$ [nats]
           & $I(\text{interv.}; s)$ [bits]
           & $I(\text{target}; s)$ [bits]
           & $\eta$
           & posts/episode \\
    \midrule
    0.75 & $\mathbf{3.045}$ \,[2.77, 3.28] & $\mathbf{0.294}$ \,[0.26, 0.33] & 0.022 & 0.006 & 69.1 \\
    1.00 & $\mathbf{1.931}$ \,[1.60, 2.31] & $\mathbf{0.154}$ \,[0.12, 0.19] & 0.004 & 0.001 & 83.9 \\
    \bottomrule
  \end{tabular}
  \caption{Bias-corrected cluster-bootstrap intervals (300 resamples, clustered
  by initialization). Target information and $\eta$ are \emph{confounded} by the
  pool defect and are reported for completeness only.}
  \label{tab:estimates}
\end{table}

\begin{figure}[t]
  \centering
  \includegraphics[width=\textwidth]{fig_pathkl}
  \caption{Path divergence from silence as a function of horizon. Intervention is
  \emph{more} displacing at $\rho = 0.75$ than at $\rho = 1$ despite fewer posts
  per episode, consistent with a refresh mechanism: at finite persistence the
  controller's facts keep being re-injected into decaying memories, whereas under
  perfect memory one post is enough and further posts are redundant.}
  \label{fig:pathkl}
\end{figure}

\begin{figure}[t]
  \centering
  \includegraphics[width=\textwidth]{fig_mi}
  \caption{Left: intervention is clearly visible in the coarse state ($0.29$ bits
  of a possible $1$ at $\rho = 0.75$). Right: the \emph{direction} of steering is
  not --- note the axis scale. This is the quantitative signature of the pool
  defect: both arms quote the same facts, so the trajectory cannot tell you which
  target was configured.}
  \label{fig:mi}
\end{figure}

\begin{figure}[t]
  \centering
  \includegraphics[width=0.62\textwidth]{fig_bins}
  \caption{Choosing the number of bins. Held-out order-2 gain (lower is more
  Markovian) against bin count, for quantile and semantic cut points. Dotted line
  marks the chosen value. Differences among 3, 4 and 5 bins do not survive
  resampling across seeds.}
  \label{fig:bins}
\end{figure}

\subsection{Three warnings about these numbers}

\paragraph{Validity is cell-specific.} Only $\rho = 1$ / silent passes all four
Markov tests. At $\rho = 0.75$ the numbers carry about $0.015$ nats/step of
unexplained memory and a Chapman--Kolmogorov gap reaching $0.11$ at five steps.
Long-horizon path quantities there are indicative, not exact.

\paragraph{\texttt{EP} is not physical dissipation.} The column we compute is the
Markov-chain irreversibility
$\sum_{ij} p_i T_{ij} \ln (p_i T_{ij} / p_j T_{ji})$. No reverse protocol has
been justified for this system, and a sequence-reversal score must not be
relabelled entropy production. It also \emph{decreases} with horizon here, which
is what one expects of a transient chain rather than a stationary one.

\paragraph{The bootstrap needed a bias correction.} Resampling episodes with
replacement duplicates them, which smooths the refitted kernel and pulls any
divergence systematically toward zero. Raw percentile intervals therefore
excluded their own point estimates. The intervals in
Table~\ref{tab:estimates} are shifted by (bootstrap mean $-$ point estimate).

\section{What this means for the next experiment}

\begin{enumerate}
  \item \textbf{Re-run with matched symmetric pools.} Target information cannot
        be interpreted until the two arms quote different, comparably strong
        evidence.
  \item \textbf{Keep $\rho = 1$ as the primary cell.} It is the only one where
        the path divergence is exact rather than approximate. Finite persistence
        remains interesting --- it is where intervention bites hardest --- but
        needs either the 8-state coordinate or substantially more episodes.
  \item \textbf{Decide the controller's role deliberately.} It currently holds a
        pool that can prove nothing (0 of 7{,}694 minimal proofs), while the
        agents jointly hold 4. Whether the controller should be able to prove its
        target is a design choice, not a detail, and it determines whether the
        experiment measures persuasion or disclosure.
  \item \textbf{Budget for the estimator, not just the simulation.} The binding
        constraint on every quantity in Table~\ref{tab:estimates} is transitions
        per (arm, $\rho$) cell, not episodes.
\end{enumerate}

\section*{Reproducing this}

\begin{center}
\small
\begin{tabular}{ll}
\toprule
step & command \\
\midrule
per-round states & \texttt{scripts/build\_states.py} \\
epistemic coordinate & \texttt{scripts/add\_coordinates.py} \\
the four tests & \texttt{scripts/run\_four\_tests.py} \\
bin sweep & \texttt{scripts/bin\_sweep.py} \\
estimates & \texttt{scripts/estimate.py} \\
this document & \texttt{scripts/build\_report.py} \\
\bottomrule
\end{tabular}
\end{center}

\noindent
Inputs: the archived study under
\texttt{agents\_control/new\_rnd\_init\_experiment} (read-only) and the generator
source in \texttt{src/mas\_cc/musr\_team\_allocation\_generator}. This analysis lives under
\texttt{analysis/preanalysis\_task003\_new\_setup\_and\_coarsegraining/} and is versioned on
branch \texttt{dev/rsanchez}; its \texttt{results/} subdirectory (per-round state
files, logs) is not.

\end{document}
